% !PRINTED 50
\weberpage{50}
%<<CONT
Gruppe der $B$, die also auch ein Theiler von $S$ ist, wollen wir
mit $U$ bezeichnen und die zu $T$ \gesperrt{reciproke Gruppe} nennen. Auch
hier ist aber zu bemerken, dass der Begriff der reciproken Gruppe
im Allgemeinen von der Wahl der Basis und der Einheitswurzeln
$\varepsilon$ abhängt, also nicht zu den Gruppen $S$ und $T$ als
solchen gehört.

Die zu $T'$ reciproke Gruppe ist durch folgenden Satz charakterisirt:
% [sic: Weber prints T' here; the newly defined group is reciprocal to T.]

8.\ \gesperrt{Lässt man $A$ die Elemente eines Theilers $T$ von $S$
durchlaufen und sucht alle Elemente $B$ von $S$, die für jedes $A$
der Bedingung}
\weq{7}{\chi_B(A) = 1}
\gesperrt{genügen, so durchläuft $B$ die Elemente der zu $T$
reciproken Gruppe $U$, deren Grad gleich dem Index von $T$ ist.}

Nach dem Satze \S.\,11, (20) ist $T'$ die reciproke Gruppe zu $U$,
% [sic: Weber prints T' here; reciprocity requires T.]
die Beziehung dieser beiden Gruppen also eine gegenseitige.

Von diesen reciproken Gruppen gilt noch der Satz:

9.\ \gesperrt{Ist $T'$ ein Theiler von $T$, so ist umgekehrt die
reciproke Gruppe $U$ von $T'$ ein Theiler der zu $T''$ reciproken
Gruppe $U'$.}
% [sic: Weber prints T' and T'' here; the theorem requires T and T', respectively.]

Denn jedes Element $A'$ von $T''$ ist zugleich in $T'$ enthalten,
% [sic: Weber prints T'' and T' here; the proof requires T' and T, respectively.]
und folglich ist, wenn $B$ die Gruppe $U$ durchläuft,
$\chi_B(A') = 1$. Es muss also $B$ in der Gruppe $U'$ vorkommen.

Wählt man aus der Gesammtheit der Charaktere $\chi$ eine beliebige
Anzahl $\xi_1, \xi_2 \dots \xi_\mu$ aus, gleichviel ob sie eine
Gruppe bilden oder nicht, so bilden alle Elemente $A$, die den
$\mu$ Bedingungen
\weq{8}{\xi_1(A) = 1,\ \xi_2(A) = 1, \dots \xi_\mu(A) = 1}
genügen, eine Gruppe, weil aus $\xi_i(A) = 1$, $\xi_i(A') = 1$
folgt, dass auch $\xi_i(AA') = 1$ ist. Wenn die
$\xi_1, \xi_2 \dots \xi_\mu$ keine Gruppen sind, so folgen aus den
Gleichungen (8) noch so viele weitere $\xi_{\mu+1}(A) = 1 \dots$,
dass die $\xi_1, \xi_2 \dots \xi_\mu$ zu einer Gruppe ergänzt
werden.

10.\ \gesperrt{Aus dem Theorem 7. folgt, dass man so alle
möglichen Theiler von $S$ erhalten kann.}
